%% kapa: exam=1 round=28 year=2017 period=2 form=A subject=appraisal_law subject_ko=감정평가관계법규 seq=32 number=32
%% source: Q-Net 감정평가사 기출문제 (cst00309) artlSeq=5199603 file=extracted/1008/5199603/2017년 제28회 감정평가사 1차시험 문제지/2교시 A형.pdf
%% conversion: kordoc 4.24.51 -> tools/gen1.py (auto); review=proofread-sonnet; flags=none
\begin{question}{32}
부동산등기법령상 등기의 순위와 접수 등에 관한 설명으로 옳지 않은 것은?
\begin{choices}
  \item 같은 부동산에 관하여 등기한 권리의 순위는 법률에 다른 규정이 없으면 등기한 순서에 따른다.
  \item 등기의 순서는 등기기록 중 같은 구(區)에서 한 등기 상호간에는 순위번호에 따른다.
  \item 같은 주등기에 관한 부기등기 상호간의 순위는 그 등기 순서에 따른다.
  \item 등기신청은 대법원규칙으로 정하는 등기신청정보가 전산정보처리조직에 저장된 때 접수된 것으로 본다.
  \item 등기관이 등기를 마쳤을 때에는 신청인에게 그 사실을 알려야 하며, 신청인이 등기완료의 통지를 받은 때부터 그 등기의 효력이 발생한다.
\end{choices}
\end{question}
\answer{5}
